% Unit 117 English translation of chapter8.tex lines 1704--1849.
% Source: Methods of Algebra, Volume 2, commit
% 9a5803ff77dd3257484cb177f851a73770a59dd3 (CC BY 4.0).
% stable-unit-id: o014.aljabr2.chapter8.closed-structure
% Coverage: all of Section 8.9, Closed Structures, ending immediately before
% Section 8.10.
% Mathematical/editorial correction: the opening distinguishes the Cartesian
% closed structure on sSet from the closed monoidal structure on sAb; the
% direction of f^* on source line 1710 is reversed to agree with precomposition
% and with lines 1717--1722.

% segment-id: o014.aljabr2.chapter8.closed-structure.h001
\section{Closed Structures}\label{sec:closed-simplicial}
\hypertarget{o014.aljabr2.chapter8.closed-structure}{ }
% segment-id: o014.aljabr2.chapter8.closed-structure.p001
The aim of this section is to explain the closed structures on $\cate{sSet}$
and $\cate{sAb}$. For $\cate{sSet}$, this is the Cartesian closed structure
of Definition~\ref{def:Cartesian-closed}; for $\cate{sAb}$, it is the
monoidal structure given by the tensor product that is closed. In other
words, we seek to promote the $\Hom$ sets to suitable simplicial objects. We
first discuss the category $\cate{sSet}$ of simplicial sets, then its version
for $\cate{sAb}$ and its comparison, via the Dold--Kan correspondence, with
the $\Hom$ complex. The last step requires the results of
\S\ref{sec:bisimplicial}.

% segment-id: o014.aljabr2.chapter8.closed-structure.d001
\begin{definition}\label{def:internal-Hom-sSet}
	For $X,Y\in\Obj(\cate{sSet})$, define
	$\iHom(X,Y)=\iHom_{\cate{sSet}}(X,Y)\in\Obj(\cate{sSet})$ as follows:
	% segment-id: o014.aljabr2.chapter8.closed-structure.q001
	\[
	\iHom(X,Y)_n:=\Hom_{\cate{sSet}}(X\times\Delta^n,Y),
	\qquad n\in\Z_{\geq0}.
	\]
	For every morphism $f:[m]\to[n]$ in $\simpDelta$, define
	$f^*:\iHom(X,Y)_n\to\iHom(X,Y)_m$ by pullback, that is, by
	precomposition, along
	% segment-id: o014.aljabr2.chapter8.closed-structure.q002
	\[
	\identity_X\times f:X\times\Delta^m\to X\times\Delta^n.
	\]
	Also define the evaluation morphism
	% segment-id: o014.aljabr2.chapter8.closed-structure.q003
	\[
	\mathrm{ev}_{X,Y}:\iHom(X,Y)\times X\to Y
	\]
	as follows: the image of
	$(\varphi,x)\in\Hom_{\cate{sSet}}(X\times\Delta^n,Y)\times X_n$ is
	$\mathrm{ev}_{X,Y,n}(\varphi,x):=\varphi(x,\identity_{[n]})\in Y_n$.
	Recall that $(\Delta^n)_n=\End_{\simpDelta}([n])$.
\end{definition}

% segment-id: o014.aljabr2.chapter8.closed-structure.p002
We must check that $\mathrm{ev}_{X,Y}$ is indeed a morphism in
$\cate{sSet}$. This is immediate: for $f:[m]\to[n]$ and
$(\varphi,x)\in\iHom(X,Y)_n\times X_n$, we have
% segment-id: o014.aljabr2.chapter8.closed-structure.q004
\begin{align*}
	\mathrm{ev}_{X,Y,m}\bigl(f^*(\varphi),f^*(x)\bigr)
	&=((\identity_X\times f)^*\varphi)
	  \bigl(f^*(x),\identity_{[m]}\bigr)\\
	&=\varphi\bigl(
	  \underbracket{f^*(x),\;f}_{\in X_m\times(\Delta^n)_m}\bigr)
	 =\varphi\bigl(f^*(x),f^*\identity_{[n]}\bigr)\\
	&=f^*\bigl(\varphi(x,\identity_{[n]})\bigr)
	 =f^*\bigl(\mathrm{ev}_{X,Y,n}(\varphi,x)\bigr).
\end{align*}

% segment-id: o014.aljabr2.chapter8.closed-structure.p003
As $X,Y$ vary, this construction gives a functor
$\iHom:\cate{sSet}^{\opp}\times\cate{sSet}\to\cate{sSet}$, while
$\mathrm{ev}$ is natural in $X$ and $Y$.

% segment-id: o014.aljabr2.chapter8.closed-structure.pr001
\begin{proposition}\label{prop:sSet-closed}
	For every $X,Y,Z\in\Obj(\cate{sSet})$, there is a canonical bijection
	% segment-id: o014.aljabr2.chapter8.closed-structure.g001
	\[\begin{tikzcd}
		\Hom_{\cate{sSet}}\bigl(X,\iHom(Y,Z)\bigr)
		\arrow[r,"1:1"]
		&\Hom_{\cate{sSet}}(X\times Y,Z).
	\end{tikzcd}\]
	It sends $g:X\to\iHom(Y,Z)$ to the composite
	% segment-id: o014.aljabr2.chapter8.closed-structure.q005
	\[
	X\times Y\xrightarrow{g\times\identity_Y}
	\iHom(Y,Z)\times Y\xrightarrow{\mathrm{ev}_{Y,Z}}Z.
	\]
	Conversely, it sends $h:X\times Y\to Z$ to the following morphism. For
	$n\in\Z_{\geq0}$ and $x\in X_n$, let $\iota_x:\Delta^n\to X$ be the
	morphism corresponding to $x$. The image of $x$ is the composite
	% segment-id: o014.aljabr2.chapter8.closed-structure.q006
	\[
	Y\times\Delta^n\xrightarrow{\identity_Y\times\iota_x}Y\times X
	\xrightarrow[\sim]{\text{swap}}X\times Y\xrightarrow{h}Z.
	\]
\end{proposition}
% segment-id: o014.aljabr2.chapter8.closed-structure.pf001
\begin{proof}
	A routine verification.
\end{proof}

% segment-id: o014.aljabr2.chapter8.closed-structure.c001
\begin{corollary}
	The bijection in Proposition~\ref{prop:sSet-closed} makes $\cate{sSet}$
	a Cartesian closed category in the sense of
	Definition~\ref{def:Cartesian-closed}, with bifunctor $\iHom$.
\end{corollary}
% segment-id: o014.aljabr2.chapter8.closed-structure.pf002
\begin{proof}
	Recall that the product in $\cate{sSet}$ is simply the product
	$X\times Y$ of simplicial sets, and similarly for the other products.
	It remains to compare \eqref{eqn:closed-monoidal-adj0} with the canonical
	bijection in Proposition~\ref{prop:sSet-closed}.
\end{proof}

% segment-id: o014.aljabr2.chapter8.closed-structure.p004
The general theory of closed monoidal categories (see
\S\ref{sec:closed-monoidal}) gives a canonical morphism
$\iHom(X,Y)\times X\to Y$: it is the image of $\identity$ under
% segment-id: o014.aljabr2.chapter8.closed-structure.q007
\[
\End_{\cate{sSet}}\bigl(\iHom(X,Y)\bigr)\xrightarrow{1:1}
\Hom_{\cate{sSet}}\bigl(\iHom(X,Y)\times X,Y\bigr).
\]
Unwinding the definitions shows that this morphism is precisely the
$\mathrm{ev}_{X,Y}$ defined above.

% segment-id: o014.aljabr2.chapter8.closed-structure.eg001
\begin{example}
	Proposition~\ref{prop:nerve-ff} says that the nerve functor $\mathrm{N}$
	embeds $\cate{Cat}$ fully faithfully into $\cate{sSet}$. How is $\iHom$
	reflected at the level of categories? The answer is simple: for any small
	categories $\mathcal{C}$ and $\mathcal{D}$, there is a canonical
	isomorphism
	% segment-id: o014.aljabr2.chapter8.closed-structure.q008
	\[
	\iHom(\mathrm{N}\mathcal{C},\mathrm{N}\mathcal{D})
	\simeq\mathrm{N}(\mathcal{D}^{\mathcal{C}}).
	\]
	Indeed, $\mathrm{N}$ takes products of categories to products of
	simplicial sets, and $\mathrm{N}([n])=\Delta^n$. Since $\cate{Cat}$ is
	also Cartesian closed, we obtain
	% segment-id: o014.aljabr2.chapter8.closed-structure.q009
	\begin{multline*}
		\Hom_{\cate{sSet}}(\mathrm{N}\mathcal{C}\times\Delta^n,
		\mathrm{N}\mathcal{D})
		=\Hom_{\cate{sSet}}(\mathrm{N}(\mathcal{C}\times[n]),
		\mathrm{N}\mathcal{D})\\
		\simeq\Hom_{\cate{Cat}}(\mathcal{C}\times[n],\mathcal{D})
		\simeq\Hom_{\cate{Cat}}([n],\mathcal{D}^{\mathcal{C}})
		=\mathrm{N}(\mathcal{D}^{\mathcal{C}})_n.
	\end{multline*}
	It is easy to check that the pullbacks on both sides correspond for every
	morphism $f:[m]\to[n]$ in $\simpDelta$.
\end{example}

% segment-id: o014.aljabr2.chapter8.closed-structure.p005
Now consider $\cate{sAb}$. Recall that it has the symmetric monoidal
structure arising from $\otimes:=\otimes_{\Z}$, namely the degreewise tensor
product. For $X\in\Obj(\cate{sAb})$ and $K\in\Obj(\cate{sSet})$, define
% segment-id: o014.aljabr2.chapter8.closed-structure.q010
\[
X\otimes K:=X\otimes\Z K\in\Obj(\cate{sAb}).
\]
In particular, $X\otimes\Delta^n$ is defined. This construction gives a
bifunctor $\cate{sAb}\times\cate{sSet}\to\cate{sAb}$.

% segment-id: o014.aljabr2.chapter8.closed-structure.d002
\begin{definition}
	For $X,Y\in\Obj(\cate{sAb})$, define
	$\iHom(X,Y)=\iHom_{\cate{sAb}}(X,Y)\in\Obj(\cate{sAb})$ by
	$\iHom(X,Y)_n:=\Hom_{\cate{sAb}}(X\otimes\Delta^n,Y)$.
	The pullback morphisms are defined as in
	Definition~\ref{def:internal-Hom-sSet}.
\end{definition}

% segment-id: o014.aljabr2.chapter8.closed-structure.pr002
\begin{proposition}\label{prop:sAb-closed}
	For every $X,Y,Z\in\Obj(\cate{sAb})$, there is a canonical commutative
	diagram
	% segment-id: o014.aljabr2.chapter8.closed-structure.g002
	\[\begin{tikzcd}
		\Hom_{\cate{sSet}}\bigl(X,\iHom_{\cate{sSet}}(Y,Z)\bigr)
		\arrow[r,"1:1"]
		&\Hom_{\cate{sSet}}(X\times Y,Z)\\
		\Hom_{\cate{sAb}}\bigl(X,\iHom_{\cate{sAb}}(Y,Z)\bigr)
		\arrow[r,"1:1"']
		\arrow[phantom,u,"\subset" description,sloped]
		&\Hom_{\cate{sAb}}(X\otimes Y,Z)\arrow[hookrightarrow,u]
	\end{tikzcd}\]
\end{proposition}
% segment-id: o014.aljabr2.chapter8.closed-structure.pf003
\begin{proof}
	The first row is Proposition~\ref{prop:sSet-closed}. The vertical arrow on
	the right embeds $\Hom_{\cate{sAb}}(X\otimes Y,Z)$ as the bilinear part
	of $\Hom_{\cate{sSet}}(X\times Y,Z)$. The image of the embedding on the
	left consists of all $\varphi:X\to\iHom_{\cate{sSet}}(Y,Z)$ for which
	the family of morphisms
	% segment-id: o014.aljabr2.chapter8.closed-structure.q012
	\[
	\varphi_n:X_n\to\Hom_{\cate{sSet}}(Y\times\Delta^n,Z)
	\]
	satisfies, for every $n\in\Z_{\geq0}$,
	% segment-id: o014.aljabr2.chapter8.closed-structure.l001
	\begin{itemize}
		\item $\varphi_n(x):Y\times\Delta^n\to Z$ is additive in the $Y$
		variable for every $x\in X_n$, or equivalently factors through
		$\Hom_{\cate{sAb}}(Y\otimes\Delta^n,Z)$;
		\item the resulting map
		$\varphi_n:X_n\to\Hom_{\cate{sAb}}(Y\otimes\Delta^n,Z)$ is additive.
	\end{itemize}
	Unwinding the definitions immediately shows that these two descriptions
	correspond.
\end{proof}

% segment-id: o014.aljabr2.chapter8.closed-structure.p006
In fact, $\cate{sAb}$ is an additive category, and the bijection
$\Hom_{\cate{sAb}}\bigl(X,\iHom_{\cate{sAb}}(Y,Z)\bigr)
\simeq\Hom_{\cate{sAb}}(X\otimes Y,Z)$ is an isomorphism of $\Z$-modules.

% segment-id: o014.aljabr2.chapter8.closed-structure.c002
\begin{corollary}
	Let $X,Y\in\Obj(\cate{sAb})$ and $K\in\Obj(\cate{sSet})$. There are
	canonical isomorphisms of $\Z$-modules
	% segment-id: o014.aljabr2.chapter8.closed-structure.q013
	\begin{align*}
		\Hom_{\cate{sAb}}(X\otimes K,Y)
		&\simeq\Hom_{\cate{sSet}}
		\bigl(K,\iHom_{\cate{sAb}}(X,Y)\bigr)\\
		&\simeq\Hom_{\cate{sAb}}
		\bigl(X,\iHom_{\cate{sSet}}(K,Y)\bigr).
	\end{align*}
	By the universal property, the underlying simplicial set of
	$\iHom_{\cate{sSet}}(K,Y)$ in the last term is naturally identified with
	the underlying simplicial set of $\iHom_{\cate{sAb}}(\Z K,Y)$, and thus
	acquires the structure of an object of $\cate{sAb}$.
\end{corollary}
% segment-id: o014.aljabr2.chapter8.closed-structure.pf004
\begin{proof}
	Proposition~\ref{prop:sAb-closed} gives
	% segment-id: o014.aljabr2.chapter8.closed-structure.q014
	\begin{align*}
		\Hom_{\cate{sAb}}(X\otimes K,Y)
		&\simeq\Hom_{\cate{sAb}}(\Z K\otimes X,Y)\\
		&\simeq\Hom_{\cate{sAb}}
		\bigl(\Z K,\iHom_{\cate{sAb}}(X,Y)\bigr).
	\end{align*}
	The adjunction property of $\Z(\cdot)$ shows that the last term is
	$\Hom_{\cate{sSet}}(K,\iHom_{\cate{sAb}}(X,Y))$. Similarly,
	$\Hom_{\cate{sAb}}(X\otimes K,Y)\simeq
	\Hom_{\cate{sAb}}(X,\iHom_{\cate{sAb}}(\Z K,Y))$.
\end{proof}

% segment-id: o014.aljabr2.chapter8.closed-structure.c003
\begin{corollary}
	The canonical bijection in Proposition~\ref{prop:sAb-closed} makes
	$(\cate{sAb},\otimes)$ a closed monoidal category in the sense of
	Definition~\ref{def:closed-monoidal-cat}, with bifunctor $\iHom$.
\end{corollary}

% segment-id: o014.aljabr2.chapter8.closed-structure.p007
As for simplicial sets, the general theory of closed monoidal categories
gives a canonical evaluation morphism
% segment-id: o014.aljabr2.chapter8.closed-structure.q015
\[
\mathrm{ev}_{X,Y}:\iHom_{\cate{sAb}}(X,Y)\otimes X\to Y.
\]
This morphism has a concrete description similar to
Definition~\ref{def:internal-Hom-sSet}, obtained simply by unwinding the
definitions.

% segment-id: o014.aljabr2.chapter8.closed-structure.p008
We now study the relationship between $\iHom_{\cate{sAb}}$ and the $\Hom$
complex from the perspective of the Dold--Kan correspondence
(Theorem~\ref{prop:Dold-Kan}). They are the internal $\Hom$ objects of the
closed monoidal categories $\cate{sAb}$ and
$\cate{Ch}_{\geq0}(\cate{Ab})$, respectively. Although the normalized chain
complex functor
$\mathrm{N}:\cate{sAb}\to\cate{Ch}_{\geq0}(\cate{Ab})$ is an equivalence,
it is not a monoidal functor. We therefore cannot directly deduce an
isomorphism between $\mathrm{N}\iHom_{\cate{sAb}}(X,Y)$ and
$\Hom^\bullet(\mathrm{N}X,\mathrm{N}Y)$. Nevertheless, the results of
\S\ref{sec:bisimplicial} still provide a pair of canonical morphisms
% segment-id: o014.aljabr2.chapter8.closed-structure.g003
\begin{equation}\label{eqn:Hom-N}
\begin{tikzcd}
	\mathrm{N}\iHom_{\cate{sAb}}(X,Y)\arrow[shift left,r]
	&\Hom^\bullet(\mathrm{N}X,\mathrm{N}Y).\arrow[shift left,l]
\end{tikzcd}
\end{equation}
% segment-id: o014.aljabr2.chapter8.closed-structure.l002
\begin{itemize}
	\item By the closed monoidal structure on
	$\cate{Ch}_{\geq0}(\cate{Ab})$, giving a morphism from left to right is
	equivalent to giving a morphism
	% segment-id: o014.aljabr2.chapter8.closed-structure.q016
	\[
	\mathrm{N}\iHom_{\cate{sAb}}(X,Y)\otimes\mathrm{N}X
	\longrightarrow\mathrm{N}Y.
	\]
	Take this morphism to be the composite
	% segment-id: o014.aljabr2.chapter8.closed-structure.q017
	\[
	\mathrm{N}\iHom_{\cate{sAb}}(X,Y)\otimes\mathrm{N}X
	\xrightarrow{\mathrm{EZ}}
	\mathrm{N}\bigl(\iHom_{\cate{sAb}}(X,Y)\otimes X\bigr)
	\xrightarrow{\mathrm{N}(\mathrm{ev}_{X,Y})}\mathrm{N}Y.
	\]
	\item Apply the left adjoint $\Gamma$ of $\mathrm{N}$, which is also its
	quasi-inverse. Giving a morphism from right to left is equivalent to giving
	a morphism
	% segment-id: o014.aljabr2.chapter8.closed-structure.q018
	\[
	\Gamma\Hom^\bullet(\mathrm{N}X,\mathrm{N}Y)
	\longrightarrow\iHom_{\cate{sAb}}(X,Y).
	\]
	By the closed monoidal structure on $\cate{sAb}$, this is also equivalent
	to giving a morphism
	$\Gamma\Hom^\bullet(\mathrm{N}X,\mathrm{N}Y)\otimes X\to Y$.
	Concretely, choose this morphism so that its image under $\mathrm{N}$ is
	the composite
	% segment-id: o014.aljabr2.chapter8.closed-structure.q019
	\begin{multline*}
		\mathrm{N}\bigl(\Gamma\Hom^\bullet(\mathrm{N}X,\mathrm{N}Y)
		\otimes X\bigr)
		\xrightarrow{\mathrm{AW}}
		\mathrm{N}\Gamma\Hom^\bullet(\mathrm{N}X,\mathrm{N}Y)
		\otimes\mathrm{N}X\\
		\simeq\Hom^\bullet(\mathrm{N}X,\mathrm{N}Y)\otimes\mathrm{N}X
		\xrightarrow{\text{evaluation}}\mathrm{N}Y.
	\end{multline*}
\end{itemize}

% segment-id: o014.aljabr2.chapter8.closed-structure.p009
The point of the construction above is that $\mathrm{N}$ is both a right
lax and a left lax monoidal functor. If $\mathrm{N}$ were monoidal, or
equivalently if $\mathrm{EZ}$ and $\mathrm{AW}$ were both isomorphisms, an
abstract argument would show that the two morphisms in \eqref{eqn:Hom-N}
are mutually inverse; in reality, they are not. However, the general
Eilenberg--Zilber theorem, Theorem~\ref{prop:Eilenberg-Zilber}, states that
$\mathrm{EZ}$ and $\mathrm{AW}$ are mutually inverse on homology. This gives
the following relationship between the $\Hom$ complex and
$\iHom_{\cate{sAb}}$.

% segment-id: o014.aljabr2.chapter8.closed-structure.pr003
\begin{proposition}
	For every $X,Y\in\Obj(\cate{sAb})$, the pair of morphisms in
	\eqref{eqn:Hom-N} are mutually inverse on homology.
\end{proposition}

% segment-id: o014.aljabr2.chapter8.closed-structure.p010
If $\Bbbk$ is a commutative ring, an analogous result holds after replacing
$\cate{Ab}$ by the symmetric monoidal category $\Bbbk\dcate{Mod}$ and
considering simplicial objects in it.

% segment-id: o014.aljabr2.chapter8.closed-structure.r001
\begin{remark}
	The definition of $\iHom$ can be extended to an arbitrary additive
	category $\mathcal{A}$. Indeed, in the description of
	$\iHom_{\cate{sAb}}(X,Y)_n$, we have
	% segment-id: o014.aljabr2.chapter8.closed-structure.q020
	\[
	(X\otimes\Z\Delta^n)_k
	=X_k\otimes\left(\bigoplus_{u:[k]\to[n]}\Z\right)
	\simeq\bigoplus_{u:[k]\to[n]}X_k.
	\]
	The last term involves only finite direct sums and no elements; the maps
	$d_i,s_j$, and so forth also have simple descriptions in this direct-sum
	presentation. We can therefore define a bifunctor
	% segment-id: o014.aljabr2.chapter8.closed-structure.q021
	\[
	\iHom_{\mathcal{A}}:(\cate{s}\mathcal{A})^{\opp}
	\times\cate{s}\mathcal{A}\to\cate{s}\mathcal{A},
	\]
	which is additive in each variable.
\end{remark}
